\chapter{Intrinsic Geometry and the Algebra of Tensors}

Tensor notation is useful only when it does some of the thinking for us.
Well-designed notation behaves like a box of mechanical pieces: an upper
index fits a lower index, unlike objects refuse to be added, and a legal
contraction leaves behind an object whose transformation law can be read
from the remaining indices. The immediate purpose of this machinery is
geometric. We want to distinguish genuine curvature from curvature that
appears only because the coordinates have been bent.

\section{Curvature Seen from Inside}

Begin with an ordinary sheet of paper. It may be rolled into a cylinder
without stretching or tearing. To an observer in the surrounding room,
the sheet is plainly bent; nevertheless, distances measured within the
sheet are unchanged. Its extrinsic curvature has changed, while its
intrinsic geometry has not.

Imagine now a small creature confined to the sheet. It cannot look into
the ambient room. It can measure lengths, compare neighboring directions,
construct triangles, and transport measuring rods from place to place.
Those internal experiments are all that geometry should require. The
corresponding question in gravitation is whether a field represents true
spacetime curvature or merely a curvilinear description of flat
spacetime.

One way to make this intrinsic viewpoint concrete is to replace the
surface by a network of points and links. Join neighboring points,
record the length of every link, and allow the joints to hinge. An
equilateral triangular lattice lies flat. Bending the hinges does not
alter the encoded geometry because the lattice can always be returned
to the plane without changing a length.

Now lengthen the six links meeting at one vertex while holding the
surrounding links fixed. The central vertex can no longer remain in the
plane without violating the prescribed lengths. It must rise into a
bulge. No view from outside was needed to diagnose this curvature: the
incompatible neighboring distances already contained it.

\begin{figure}[H]
\centering
\includegraphics[width=\linewidth]{lecture_02_figure_01.png}
\caption{The triangular-link model of intrinsic geometry. Altering the
links around one vertex forces a bulge that cannot be flattened without
changing lengths.}
\lectureframe{00:06:40}
\end{figure}

\begin{classroomqa}
\audiencequestion{What, then, is the precise definition of a flat
space?}
\lecturerresponse{The operational answer is that its intrinsic distance
relations can be realized without distortion in a Euclidean layout. In
the language developed below, a Riemannian space is flat when one can
find coordinates, throughout the region of interest, in which the
metric is the constant Kronecker delta. At a single point this can
always be done locally; doing it over a region is the nontrivial
condition. The curvature tensor will eventually state that condition
without reference to a special coordinate choice.
\lecturetimestamp{00:09:15}}
\end{classroomqa}

This answer uses the metric before we have earned it. We therefore turn
temporarily from curvature to tensor analysis. The detour is not
optional: the metric must be recognized as a tensor before its
coordinate dependence can be separated from genuine geometry.

\section{Fields and Coordinate Descriptions}

A tensor in general relativity is usually a tensor \emph{field}: at
every point there is a tensor, and its components vary from point to
point. A scalar field is the simplest example. Temperature on a surface
may be represented by a number \(S(P)\) at each point \(P\). Two people
may label the point by different coordinates, but they must agree on
the temperature there.

\begin{figure}[H]
\centering
\includegraphics[width=\linewidth]{lecture_02_figure_02.png}
\caption{A scalar field together with two coordinate labels for the
same points.}
\lectureframe{00:11:25}
\end{figure}

Let \(x^m\) and \(y^m\) be two smooth coordinate systems. Each is a map
from points to ordered lists of numbers, and in an overlapping
coordinate patch the two lists are related by mutually inverse
functions,
\begin{equation}
 y^m=y^m(x^1,\ldots,x^d),
 \qquad
 x^m=x^m(y^1,\ldots,y^d).
\label{eq:coordinate-maps}
\end{equation}
Smoothness allows differentiation, and invertibility ensures that no
point is lost or identified with another merely by changing
coordinates.

\begin{figure}[H]
\centering
\includegraphics[width=\linewidth]{lecture_02_figure_03.png}
\caption{Two smooth, invertible coordinate maps describing the same
geometric region.}
\lectureframe{00:13:10}
\end{figure}

For a scalar field the transformation law is simply
\begin{equation}
 S'(y)=S(x).
\label{eq:scalar-law}
\end{equation}
The functions on the two sides may look different because their
arguments differ, but at corresponding points their values agree. This
is the prototype of an invariant statement.

\section{Oblique Coordinates and Two Kinds of Components}

Before defining general tensors by transformation laws, it is useful to
see why upper and lower vector indices cannot be treated as decorative
variants of one another.

Draw a coordinate grid whose axes are neither perpendicular nor
normalized. At a point, increasing \(x^m\) by one coordinate unit
produces a basis vector \(\mathbf e_m\). In coordinate language,
\begin{equation}
 \mathbf e_m=\frac{\partial \mathbf r}{\partial x^m}.
\label{eq:coordinate-basis}
\end{equation}
The vectors \(\mathbf e_m\) need not have unit length and need not be
orthogonal.

\begin{figure}[H]
\centering
\includegraphics[width=\linewidth]{lecture_02_figure_04.png}
\caption{An oblique coordinate grid and its coordinate basis vectors.}
\lectureframe{00:17:40}
\end{figure}

Every vector \(\mathbf V\) can be expanded in this basis:
\begin{equation}
 \mathbf V=V^m\mathbf e_m.
\label{eq:contravariant-expansion}
\end{equation}
The numbers \(V^m\) are the \emph{contravariant components}. They are
the coefficients needed to assemble the geometric vector from the
basis vectors.

\begin{figure}[H]
\centering
\includegraphics[width=\linewidth]{lecture_02_figure_05.png}
\caption{The expansion of a vector in a nonorthogonal coordinate
basis.}
\lectureframe{00:19:20}
\end{figure}

There is another natural set of numbers:
\begin{equation}
 V_m=\mathbf V\mathbin{\boldsymbol{\cdot}}\mathbf e_m.
\label{eq:covariant-projection}
\end{equation}
These are the \emph{covariant components}. In an oblique basis they are
not simply the same coefficients with the index moved. Substituting
\eqref{eq:contravariant-expansion} gives
\begin{equation}
 V_m
 =V^n\bigl(\mathbf e_n\mathbin{\boldsymbol{\cdot}}\mathbf e_m\bigr).
\label{eq:covariant-from-basis}
\end{equation}

\begin{figure}[H]
\centering
\includegraphics[width=\linewidth]{lecture_02_figure_06.png}
\caption{Covariant components obtained by taking scalar products with
the coordinate basis.}
\lectureframe{00:21:10}
\end{figure}

The scalar products of the basis vectors form the Gram matrix
\begin{equation}
 g_{mn}=\mathbf e_m\mathbin{\boldsymbol{\cdot}}\mathbf e_n.
\label{eq:metric-gram}
\end{equation}
Consequently,
\begin{equation}
 V_m=g_{mn}V^n.
\label{eq:lower-vector}
\end{equation}
This is the first appearance of the metric. It converts expansion
coefficients into projections because it records both the lengths and
mutual angles of the coordinate basis.

\begin{figure}[H]
\centering
\includegraphics[width=\linewidth]{lecture_02_figure_07.png}
\caption{The metric as the Gram matrix of the basis vectors.}
\lectureframe{00:23:40}
\end{figure}

The same matrix computes length:
\begin{align}
 \mathbf V\mathbin{\boldsymbol{\cdot}}\mathbf V
 &=(V^m\mathbf e_m)\mathbin{\boldsymbol{\cdot}}
   (V^n\mathbf e_n) \notag\\
 &=g_{mn}V^mV^n.
\label{eq:vector-length}
\end{align}

\begin{figure}[H]
\centering
\includegraphics[width=\linewidth]{lecture_02_figure_08.png}
\caption{The vector norm written in terms of the metric and
contravariant components.}
\lectureframe{00:25:00}
\end{figure}

\begin{classroomqa}
\audiencequestion{Why introduce this geometric construction before the
formal transformation laws?}
\lecturerresponse{It gives the indices a concrete meaning. Upper
components are expansion coefficients, lower components are linear
measurements of the vector, and the metric connects the two. The formal
definition will then tell us that these roles survive arbitrary smooth
coordinate changes. \lecturetimestamp{00:26:07}}
\end{classroomqa}

\begin{classroomqa}
\audiencequestion{When do the covariant and contravariant components
have the same numerical values?}
\lecturerresponse{In an orthonormal Euclidean basis,
\(g_{mn}=\delta_{mn}\), so \(V_m=V^m\) component by component. Mere
orthogonality is not enough: if the basis vectors have nonunit lengths,
the diagonal metric still supplies scale factors.
\lecturetimestamp{00:26:19}}
\end{classroomqa}

Polar coordinates provide the clean warning. Radial and angular
coordinate lines meet orthogonally, yet a fixed change
\(d\theta\) corresponds to physical distance \(r\,d\theta\). In the
plane,
\begin{equation}
 ds^2=dr^2+r^2d\theta^2,
\label{eq:polar-metric}
\end{equation}
so the basis is orthogonal but not Cartesian.

\begin{figure}[H]
\centering
\includegraphics[width=\linewidth]{lecture_02_figure_09.png}
\caption{Polar coordinates: orthogonal coordinate directions with an
angular scale factor that grows with radius.}
\lectureframe{00:28:10}
\end{figure}

\begin{classroomqa}
\audiencequestion{Are covariant vectors genuinely different vectors,
or only a second way of listing the same vector?}
\lecturerresponse{Before a metric is supplied, a vector and a covector
are different kinds of linear object. A vector belongs to the tangent
space; a covector is a linear functional on that space. A nondegenerate
metric gives a one-to-one correspondence between them, represented by
\(V_m=g_{mn}V^n\), but the correspondence should not be mistaken for
identity. \lecturetimestamp{00:29:10}}
\end{classroomqa}

\section{Transformation Laws Define the Tensor}

The safest formal definition is operational: an array is a tensor if
its components transform by the appropriate Jacobian factors.
Differentiating \eqref{eq:coordinate-maps} gives
\begin{equation}
 dy^m=\frac{\partial y^m}{\partial x^n}\,dx^n.
\label{eq:differential-law}
\end{equation}
Any contravariant vector obeys the same rule,
\begin{equation}
 V'^m(y)
 =\frac{\partial y^m}{\partial x^n}V^n(x).
\label{eq:contravariant-law}
\end{equation}

\begin{figure}[H]
\centering
\includegraphics[width=\linewidth]{lecture_02_figure_10.png}
\caption{The contravariant transformation law, modeled on the
transformation of coordinate differentials.}
\lectureframe{00:32:50}
\end{figure}

A covector transforms with the inverse Jacobian:
\begin{equation}
 W'_m(y)
 =\frac{\partial x^n}{\partial y^m}W_n(x).
\label{eq:covariant-law}
\end{equation}
The gradient of a scalar is the standard example. Applying the chain
rule to \eqref{eq:scalar-law},
\begin{equation}
 \frac{\partial S}{\partial y^m}
 =\frac{\partial x^n}{\partial y^m}
  \frac{\partial S}{\partial x^n}.
\label{eq:gradient-law}
\end{equation}

\begin{figure}[H]
\centering
\includegraphics[width=\linewidth]{lecture_02_figure_11.png}
\caption{A scalar gradient transforming as a covector by the chain
rule.}
\lectureframe{00:34:20}
\end{figure}

\begin{classroomqa}
\audiencequestion{Do upper and lower components necessarily have the
same physical units?}
\lecturerresponse{No. Coordinate components inherit the units and
scale conventions of the coordinates and their basis vectors. If one
coordinate is a length and another an angle, their components need not
even share dimensions. The complete contraction is the object whose
physical dimensions are coordinate independent.
\lecturetimestamp{00:35:41}}
\end{classroomqa}

\begin{classroomqa}
\audiencequestion{Why are the derivatives partial derivatives rather
than ordinary derivatives?}
\lecturerresponse{A field depends on several coordinates. The Jacobian
asks how one coordinate function changes when one argument varies and
the others are held fixed, so its entries are partial derivatives.
\lecturetimestamp{00:36:45}}
\end{classroomqa}

The pattern now extends mechanically. Every upper primed index brings
one factor of \(\partial y/\partial x\); every lower primed index brings
one factor of \(\partial x/\partial y\). A mixed tensor transforms as
\begin{equation}
 T'^m{}_n
 =
 \frac{\partial y^m}{\partial x^p}
 \frac{\partial x^q}{\partial y^n}
 T^p{}_q,
\label{eq:mixed-law}
\end{equation}
and a rank-two covariant tensor transforms as
\begin{equation}
 W'_{mn}
 =
 \frac{\partial x^p}{\partial y^m}
 \frac{\partial x^q}{\partial y^n}
 W_{pq}.
\label{eq:rank-two-covariant-law}
\end{equation}

\begin{figure}[H]
\centering
\includegraphics[width=\linewidth]{lecture_02_figure_12.png}
\caption{The transformation rule for a tensor with one upper and one
lower index.}
\lectureframe{00:40:45}
\end{figure}

\begin{figure}[H]
\centering
\includegraphics[width=\linewidth]{lecture_02_figure_13.png}
\caption{A rank-two covariant transformation law, with one inverse
Jacobian for each lower index.}
\lectureframe{00:42:15}
\end{figure}

Repeated indices are summed unless stated otherwise. This convention
removes summation signs while preserving all the information needed to
reconstruct them. A repeated dummy index must occur exactly twice in a
term, once upstairs and once downstairs; a free index occurs once in
every term of an equation.

\begin{classroomqa}
\audiencequestion{Who introduced this summation notation?}
\lecturerresponse{It is conventionally credited to Einstein, which is
why it is called the Einstein summation convention. Its value is not
historical ornament: the index pattern makes legal tensor operations
visible at a glance. \lecturetimestamp{00:43:21}}
\end{classroomqa}

\section{Tensor Equations and Coordinate-Independent Truth}

Components vary with coordinates, but a tensor equation remains true
in every coordinate system. If two tensors of the same type satisfy
\(W=T\), then their difference is the zero tensor:
\begin{equation}
 W^{m\cdots}{}_{\cdots n}
 -T^{m\cdots}{}_{\cdots n}=0.
\label{eq:tensor-equality}
\end{equation}
The transformation law is linear, so zero components in one system
produce zero components in every system. The matching free indices on
the two sides are indispensable. An equation between objects of
different tensor type has no coordinate-independent meaning.

\begin{figure}[H]
\centering
\includegraphics[width=\linewidth]{lecture_02_figure_14.png}
\caption{Matching index structure turns equality of components into a
tensor equation valid in every coordinate system.}
\lectureframe{00:46:35}
\end{figure}

\begin{classroomqa}
\audiencequestion{What exactly is invariant if the vector components
are changing?}
\lecturerresponse{The geometric tensor and the truth of a tensor
equation are invariant; the coordinate list is not. A scalar is the
special case with no indices, so its numerical value at a point is the
same in every coordinate system. \lecturetimestamp{00:47:19}}
\end{classroomqa}

\begin{classroomqa}
\audiencequestion{Is vector magnitude the invariant analogue, and does
zero magnitude mean the vector is zero?}
\lecturerresponse{Once a metric is given, the contraction
\(g_{mn}V^mV^n\) is a scalar and therefore an invariant squared norm.
In positive-definite geometry it vanishes only for the zero vector. In
Lorentzian spacetime a nonzero null vector can have zero norm, so the
two statements must be kept separate. \lecturetimestamp{00:48:04}}
\end{classroomqa}

\section{The Permitted Algebra of Tensors}

The board collects the principal operations that will be needed:
multiply by a scalar, add or subtract tensors of the same type, form
tensor products, contract an upper with a lower index, and eventually
differentiate covariantly.

\begin{figure}[H]
\centering
\includegraphics[width=\linewidth]{lecture_02_figure_15.png}
\caption{The board inventory of tensor operations: addition,
multiplication, contraction, and differentiation.}
\lectureframe{00:50:40}
\end{figure}

\begin{classroomqa}
\audiencequestion{Differentiation with respect to what?}
\lecturerresponse{With respect to position, hence with respect to the
coordinates. Ordinary component derivatives do not generally transform
as tensors because the Jacobians themselves vary. Covariant
differentiation repairs this, but its construction is postponed until
the geometric preliminaries are in place. \lecturetimestamp{00:51:49}}
\end{classroomqa}

Scalar multiplication is immediate. Addition and subtraction require
matching tensor type:
\begin{equation}
 C^m{}_n=A^m{}_n+B^m{}_n,
\qquad
 D^m{}_n=A^m{}_n-B^m{}_n.
\label{eq:tensor-addition}
\end{equation}
One cannot add a vector to a covector or a rank-two tensor to a scalar,
just as one cannot add a velocity to an area.

\begin{figure}[H]
\centering
\includegraphics[width=\linewidth]{lecture_02_figure_16.png}
\caption{Addition is defined componentwise only for tensors with the
same free-index pattern.}
\lectureframe{00:53:00}
\end{figure}

The tensor product combines all indices of its factors. For example,
\begin{equation}
 T^{mn}=V^mW^n,
\qquad
 S^m{}_n=V^mU_n.
\label{eq:tensor-products}
\end{equation}
In four dimensions \(T^{mn}\) has sixteen components. No index has been
summed, so this is an outer product rather than a dot product.

\begin{figure}[H]
\centering
\includegraphics[width=\linewidth]{lecture_02_figure_17.png}
\caption{A tensor product retaining one free index from each vector
factor.}
\lectureframe{00:55:15}
\end{figure}

\begin{classroomqa}
\audiencequestion{Is multiplication of two vectors here the dot
product?}
\lecturerresponse{No. A dot product contracts an index and produces a
scalar. The tensor product leaves both indices free and therefore
produces a rank-two tensor. In four dimensions that is sixteen
components rather than one. \lecturetimestamp{00:56:36}}
\end{classroomqa}

\begin{classroomqa}
\audiencequestion{Since the component numbers commute, does exchanging
the two vector factors make no difference?}
\lecturerresponse{For ordinary scalar components,
\(V^mW^n=W^nV^m\). But tensor slots are ordered. Interchanging the
factors and then naming the first slot \(m\) and the second \(n\)
produces the transposed array; \(T^{mn}\) need not equal \(T^{nm}\).
\lecturetimestamp{00:58:52}}
\end{classroomqa}

\section{Contraction and the Inverse-Jacobian Lemma}

The key to contraction is the chain rule. Because the coordinate maps
are inverse,
\begin{equation}
 \frac{\partial x^b}{\partial y^m}
 \frac{\partial y^m}{\partial x^a}
 =
 \frac{\partial x^b}{\partial x^a}
 =\delta^b{}_a.
\label{eq:inverse-jacobian}
\end{equation}
The two Jacobian matrices are inverses, and their product is the
Kronecker delta.

\begin{figure}[H]
\centering
\includegraphics[width=\linewidth]{lecture_02_figure_18.png}
\caption{The inverse-Jacobian identity prepared as the lemma behind
tensor contraction.}
\lectureframe{01:02:15}
\end{figure}

\begin{classroomqa}
\audiencequestion{What kind of object is the inverse Jacobian?}
\lecturerresponse{It is the derivative matrix of the inverse coordinate
map. Its entries are \(\partial x^b/\partial y^m\), and matrix
multiplication with \(\partial y^m/\partial x^a\) gives the identity by
the multivariable chain rule. \lecturetimestamp{01:02:37}}
\end{classroomqa}

A quick proof makes the notation transparent. For any scalar function
\(f\),
\begin{equation}
 \frac{\partial f}{\partial y^m}
 \frac{\partial y^m}{\partial x^a}
 =\frac{\partial f}{\partial x^a}.
\end{equation}
Choosing \(f=x^b\) gives
\begin{equation}
 \frac{\partial x^b}{\partial y^m}
 \frac{\partial y^m}{\partial x^a}
 =\frac{\partial x^b}{\partial x^a}
 =\delta^b{}_a.
\end{equation}

\begin{figure}[H]
\centering
\includegraphics[width=\linewidth]{lecture_02_figure_19.png}
\caption{The chain rule collapsing the inverse Jacobians to the
Kronecker delta.}
\lectureframe{01:04:55}
\end{figure}

\begin{classroomqa}
\audiencequestion{Is \(\delta^b{}_a\) the identity only in Cartesian
coordinates?}
\lecturerresponse{No. The mixed tensor \(\delta^b{}_a\) represents the
identity map in every coordinate system. A Euclidean metric written as
\(\delta_{ab}\) is special to orthonormal Cartesian coordinates, but
the mixed-index identity should not be confused with that metric.
\lecturetimestamp{01:04:58}}
\end{classroomqa}

Now take a mixed tensor \(T^m{}_n\). Its transformed trace is
\begin{align}
 T'^m{}_m
 &=
 \frac{\partial y^m}{\partial x^a}
 \frac{\partial x^b}{\partial y^m}
 T^a{}_b \notag\\
 &=\delta^b{}_aT^a{}_b
 =T^a{}_a.
\label{eq:trace-invariant}
\end{align}
The contracted quantity has no free index and is therefore a scalar.

\begin{figure}[H]
\centering
\includegraphics[width=\linewidth]{lecture_02_figure_20.png}
\caption{Transformation of a mixed tensor before the upper and lower
indices are contracted.}
\lectureframe{01:07:45}
\end{figure}

\begin{figure}[H]
\centering
\includegraphics[width=\linewidth]{lecture_02_figure_21.png}
\caption{The Jacobians cancel under contraction, leaving the same
scalar trace in either coordinate system.}
\lectureframe{01:10:30}
\end{figure}

\begin{classroomqa}
\audiencequestion{After a contraction, does the result still carry
indices?}
\lecturerresponse{Only the contracted pair disappears. Contracting the
sole upper and lower indices of \(T^m{}_n\) leaves a scalar. In a
higher-rank tensor, all uncontracted indices remain free.
\lecturetimestamp{01:09:00}}
\end{classroomqa}

For example,
\begin{equation}
 A^{mn}{}_{pq}
 \longmapsto
 A^{mn}{}_{pn}
 =B^m{}_p
\label{eq:higher-rank-contraction}
\end{equation}
contracts \(n\) and lowers the rank by two. The precise letters are
dummy labels; the structural rule is that one upper and one lower slot
are summed.

\begin{figure}[H]
\centering
\includegraphics[width=\linewidth]{lecture_02_figure_22.png}
\caption{A higher-rank contraction removes one upper and one lower
index while preserving the remaining free indices.}
\lectureframe{01:12:35}
\end{figure}

\begin{classroomqa}
\audiencequestion{Why not contract two upper indices, or two lower
indices, in the same way?}
\lecturerresponse{Two upper indices supply two Jacobians of the same
orientation,
\(\partial y^m/\partial x^a\) and
\(\partial y^m/\partial x^b\). They do not multiply to the identity, so
the proposed sum does not have a tensor transformation law. A metric
can first lower one of the indices, after which an upper--lower
contraction is legal; without that additional tensor, the operation is
not coordinate invariant. \lecturetimestamp{01:13:03}}
\end{classroomqa}

\begin{figure}[H]
\centering
\includegraphics[width=\linewidth]{lecture_02_figure_23.png}
\caption{The attempted same-variance contraction leaves two
same-direction Jacobians and therefore does not collapse to an
identity.}
\lectureframe{01:14:30}
\end{figure}

\begin{classroomqa}
\audiencequestion{Is contraction the tensor version of an inner
product?}
\lecturerresponse{A contraction pairs a vector slot with a covector
slot and is the basic invariant pairing. With a metric, two vectors can
be turned into such a pair, so the familiar inner product is a metric
followed by contraction. \lecturetimestamp{01:14:43}}
\end{classroomqa}

\begin{classroomqa}
\audiencequestion{Must the two contracted indices run over the same
range?}
\lecturerresponse{Yes. They refer to dual bases of the same
\(d\)-dimensional tangent space and cotangent space. Otherwise the
summation and the inverse-Jacobian identity would not be defined.
\lecturetimestamp{01:15:13}}
\end{classroomqa}

\begin{classroomqa}
\audiencequestion{Why does a tensor that vanishes in one frame vanish
in every frame?}
\lecturerresponse{Every transformed component is a linear combination
of the original components. All those combinations vanish when every
input component is zero. This simple linearity is what makes tensor
equations coordinate independent. \lecturetimestamp{01:15:31}}
\end{classroomqa}

\section{The Metric from Squared Distance}

Return now to geometry. For an infinitesimal displacement with
components \(dx^m\), define its squared length by
\begin{equation}
 ds^2=g_{mn}(x)\,dx^m dx^n.
\label{eq:line-element}
\end{equation}
This is the generalized Pythagorean theorem. The metric may vary from
point to point, and off-diagonal terms record nonorthogonality of the
coordinate directions.

\begin{figure}[H]
\centering
\includegraphics[width=\linewidth]{lecture_02_figure_24.png}
\caption{The tensor-operation list gives way to the quadratic
definition of the infinitesimal line element.}
\lectureframe{01:18:45}
\end{figure}

Only the symmetric part contributes because
\(dx^m dx^n=dx^n dx^m\). We may therefore take
\begin{equation}
 g_{mn}=g_{nm}.
\label{eq:metric-symmetry}
\end{equation}
In \(d\) dimensions a symmetric metric has
\begin{equation}
 \frac{d(d+1)}{2}
\label{eq:metric-component-count}
\end{equation}
independent components: ten in four dimensions, six in three, and
three in two.

\begin{figure}[H]
\centering
\includegraphics[width=\linewidth]{lecture_02_figure_25.png}
\caption{The line element and the symmetric component array used to
count the independent entries of the metric.}
\lectureframe{01:21:35}
\end{figure}

\begin{classroomqa}
\audiencequestion{Why are there ten independent metric components in
four dimensions rather than sixteen?}
\lecturerresponse{There are four diagonal entries and six unordered
off-diagonal pairs. Equivalently,
\(4(4+1)/2=10\). The antisymmetric part would multiply the symmetric
product \(dx^m dx^n\) and cancel out. \lecturetimestamp{01:19:50}}
\end{classroomqa}

\begin{classroomqa}
\audiencequestion{How many independent components are there in three
dimensions?}
\lecturerresponse{Six: three diagonal entries and three off-diagonal
pairs, or \(3(3+1)/2\). \lecturetimestamp{01:20:47}}
\end{classroomqa}

\begin{classroomqa}
\audiencequestion{And in two dimensions?}
\lecturerresponse{Three: \(g_{11}\), \(g_{22}\), and the single
symmetric pair \(g_{12}=g_{21}\). \lecturetimestamp{01:22:28}}
\end{classroomqa}

\subsection{Why the metric is a tensor}

Squared distance is a scalar. Writing the same displacement in two
coordinate systems therefore gives
\begin{equation}
 g_{mn}(x)\,dx^m dx^n
 =
 g'_{pq}(y)\,dy^pdy^q.
\label{eq:length-invariance}
\end{equation}
Use
\begin{equation}
 dx^m=\frac{\partial x^m}{\partial y^p}\,dy^p
\end{equation}
twice:
\begin{align}
 ds^2
 &=
 g_{mn}(x)
 \frac{\partial x^m}{\partial y^p}
 \frac{\partial x^n}{\partial y^q}
 dy^pdy^q.
\label{eq:metric-transform-derivation}
\end{align}
Comparison with the primed expression identifies
\begin{equation}
 g'_{pq}(y)
 =
 \frac{\partial x^m}{\partial y^p}
 \frac{\partial x^n}{\partial y^q}
 g_{mn}(x).
\label{eq:metric-transformation}
\end{equation}
This is exactly the transformation law for a covariant rank-two
tensor.

\begin{figure}[H]
\centering
\includegraphics[width=\linewidth]{lecture_02_figure_26.png}
\caption{Derivation of the metric transformation law from invariance
of the scalar line element.}
\lectureframe{01:26:25}
\end{figure}

The conclusion also explains the two lower indices. Each lower metric
index pairs with one upper index carried by a displacement component.
Under a coordinate change, the two inverse Jacobians in
\eqref{eq:metric-transformation} cancel the two forward Jacobians of
the differentials, leaving \(ds^2\) unchanged.

\section{The Metric Matrix and Its Inverse}

At a point, \(g_{mn}\) is a symmetric matrix:
\begin{equation}
 [g_{mn}]
 =
 \begin{pmatrix}
 g_{11} & g_{12} & \cdots \\
 g_{12} & g_{22} & \cdots \\
 \vdots & \vdots & \ddots
 \end{pmatrix}.
\label{eq:metric-matrix}
\end{equation}

\begin{figure}[H]
\centering
\includegraphics[width=\linewidth]{lecture_02_figure_27.png}
\caption{The metric displayed as a symmetric matrix after its tensor
transformation law has been established.}
\lectureframe{01:28:50}
\end{figure}

The metric is required to be nondegenerate: it has no zero eigenvalue
and therefore possesses an inverse. Denote that inverse by \(g^{mn}\).
In index notation,
\begin{equation}
 g_{mn}g^{np}=\delta_m{}^p.
\label{eq:inverse-metric}
\end{equation}
The repeated \(n\) is summed, while \(m\) and \(p\) label the row and
column of the product.

\begin{figure}[H]
\centering
\includegraphics[width=\linewidth]{lecture_02_figure_28.png}
\caption{The inverse metric introduced through
\(g_{mn}g^{np}=\delta_m{}^p\).}
\lectureframe{01:31:40}
\end{figure}

\begin{classroomqa}
\audiencequestion{Where do the vectors and tensors at a point actually
live?}
\lecturerresponse{Vectors at a point form its tangent space
\(T_P M\); covectors form the dual space \(T_P^*M\). Higher tensors
belong to tensor products of these spaces. A tensor field assigns such
an object smoothly to every point. The metric is a bilinear form on the
tangent space, and its inverse is the corresponding bilinear form on
the cotangent space. \lecturetimestamp{01:34:15}}
\end{classroomqa}

\begin{figure}[H]
\centering
\includegraphics[width=\linewidth]{lecture_02_figure_29.png}
\caption{A vector drawn in the tangent space attached to a point of a
curved manifold.}
\lectureframe{01:35:45}
\end{figure}

\begin{classroomqa}
\audiencequestion{Does
\(g_{mn}g^{np}=\delta_m{}^p\) refer only to the case \(m=p\)?}
\lecturerresponse{No. It is the complete matrix equation. For every
ordered pair \((m,p)\), the summed product is one when \(m=p\) and zero
when \(m\ne p\). All diagonal and off-diagonal equations are present
simultaneously. \lecturetimestamp{01:36:50}}
\end{classroomqa}

\begin{figure}[H]
\centering
\includegraphics[width=\linewidth]{lecture_02_figure_30.png}
\caption{The inverse relation interpreted as the complete
mixed-index identity matrix.}
\lectureframe{01:38:50}
\end{figure}

In two dimensions this means
\begin{equation}
 \begin{pmatrix}
 g_{11} & g_{12}\\
 g_{12} & g_{22}
 \end{pmatrix}
 \begin{pmatrix}
 g^{11} & g^{12}\\
 g^{12} & g^{22}
 \end{pmatrix}
 =
 \begin{pmatrix}
 1&0\\
 0&1
 \end{pmatrix}.
\label{eq:two-dimensional-inverse}
\end{equation}
For example,
\begin{align}
 g_{11}g^{11}+g_{12}g^{12}&=1,\\
 g_{11}g^{12}+g_{12}g^{22}&=0.
\end{align}
The other row supplies the remaining two equations.

\begin{figure}[H]
\centering
\includegraphics[width=\linewidth]{lecture_02_figure_31.png}
\caption{The inverse-metric identity expanded as an explicit
two-dimensional matrix product.}
\lectureframe{01:40:00}
\end{figure}

\begin{classroomqa}
\audiencequestion{Must the metric be positive definite?}
\lecturerresponse{For ordinary Riemannian geometry, positive
definiteness makes every nonzero squared length positive. Spacetime
uses a Lorentzian metric instead, with one sign opposite to the other
three. What both cases require here is nondegeneracy, not positivity.
\lecturetimestamp{01:42:18}}
\end{classroomqa}

\begin{classroomqa}
\audiencequestion{Why exclude a zero eigenvalue?}
\lecturerresponse{A zero eigenvalue makes the metric matrix singular,
so \(g^{mn}\) does not exist and the metric cannot establish a
one-to-one correspondence between vectors and covectors. A metric, in
the usual differential-geometric sense, is required to be
nondegenerate. \lecturetimestamp{01:43:56}}
\end{classroomqa}

\begin{classroomqa}
\audiencequestion{Do the ten independent components of a
four-dimensional metric have anything to do with ten-dimensional
string theory?}
\lecturerresponse{No. The first ten is simply
\(4(4+1)/2\), the number of entries in a symmetric four-by-four
matrix. The appearance of ten spacetime dimensions in superstring
theory has a different origin. \lecturetimestamp{01:44:32}}
\end{classroomqa}

\begin{classroomqa}
\audiencequestion{If Lorentzian spacetime has nonzero vectors of zero
length, does that make the metric noninvertible?}
\lecturerresponse{No. A null vector satisfies
\(g_{mn}V^mV^n=0\), a cancellation within an indefinite quadratic
form. It is not a zero-eigenvalue vector satisfying
\(g_{mn}V^n=0\). For example, the Minkowski metric has eigenvalues
\(\pm1\) and is perfectly invertible while still admitting lightlike
vectors. \lecturetimestamp{01:45:02}}
\end{classroomqa}

\section{What Comes Next}

We can now state the geometric program cleanly. The metric records
infinitesimal distance, transforms covariantly, and has an inverse. Its
components may look complicated merely because the coordinates are
complicated. To decide whether the geometry itself is curved, we must
compare the metric at neighboring points in a coordinate-independent
way.

Ordinary derivatives are not enough, because differentiating a tensor
also differentiates the position-dependent Jacobians in its
transformation law. The next step is therefore covariant
differentiation. From derivatives of the metric one constructs a
connection; from derivatives and products of that connection one
constructs curvature. Only then will the curled-page question receive
its final intrinsic answer.
